\section{Background}
\subsection{Directed acyclic graphs}
\begin{quote}\small\itshape
\begin{itemize}
\item DAG (directed acycclic graph)
\end{itemize}
\end{quote}
A directed acyclic graph (DAG) is a directed graph $G = (V, E)$ with no
directed cycle: there is no path that starts and ends at the same vertex
while always following edge direction. Equivalently, $G$ admits a
topological ordering of $V$ in which every edge points from an earlier to a
later vertex \citeneeded{standard graph theory reference, e.g. Cormen et al.}.
Acyclicity is what makes reachability well behaved: no vertex reachable from
$u$ by following edges forward can also reach back to $u$. Section~\ref{sec:method}
relies on this property directly -- it is what guarantees every hop-distance
used later in the paper is finite and well defined.

\subsection{SNOMED CT}
\begin{quote}\small\itshape
\begin{itemize}
\item Number of concepts, number of distinct relations,
\item DAG nature
\item semantic value from relation types and concept connectivities (concepts are defined by the others in the DAG)
\end{itemize}
\end{quote}
SNOMED~CT is a poly-hierarchy DAG containing more than 420{,}000
distinct concepts, with more than 100 distinct semantic relationships connecting concepts among
them \citeneeded{SNOMED International concept model / release statistics}.
Two properties of this graph matter. First, it is acyclic. Second, its edges
carry semantic value beyond hierarchy alone, and that value differs by
relationship type. One relationship type, \texttt{IS\_A}, encodes
generalization/specialization: an edge $u \xrightarrow{\texttt{IS\_A}} v$
means $u$ is a child of $v$, i.e.\ $u$ is strictly more specific than $v$,
matching the source/destination convention of the underlying release file
\citeneeded{SNOMED CT Release File Specification, relationship file}. Every
other relationship type instead encodes an attribute assignment rather than
a hierarchy: an edge $u \xrightarrow{R} v$, for a relation $R$ such as
\emph{finding site} or \emph{associated morphology}, means that concept
$u$'s value for attribute $R$ is the concept $v$, not that $v$ is more
general than $u$. As a consequence, a concept's meaning is not contained in
the concept itself: it is constructed from its connections inside the DAG
with respect to the other concepts. Two concepts that look unrelated by
name can share common structure if they connect via the same relationships
to the same targets. This distinction -- \texttt{IS\_A} as a hierarchy edge
versus every other relationship type as an attribute edge -- is exactly
what our method is built around: Section~\ref{sec:method} defines a
coverage score and a tree representation that treat the two kinds of edge
differently for this reason.

\subsection{Working vocabulary and notation}
With these properties established, we can now state the notation used for
the rest of the paper. Let $M \subseteq V$ denote the mapped concepts
introduced above, and let $T \subseteq V$, $|T| = k$, denote a \emph{token
set}: a small subset of concepts chosen to represent every concept in $M$.
Section~\ref{sec:method} defines exactly how $T$ is chosen and how a concept
in $M$ is represented once $T$ is fixed.